% ============================================================
%  Synthèse MTH8304 — Apprentissage non supervisé & Séries temporelles
%  Compilable sur Overleaf (pdflatex)
% ============================================================
\documentclass[10pt,a4paper,twocolumn]{article}

\usepackage[utf8]{inputenc}
\usepackage[T1]{fontenc}
\usepackage[french]{babel}
\usepackage[margin=1.2cm]{geometry}
\usepackage{amsmath,amssymb,amsthm}
\usepackage{bm}           % bold math
\usepackage{booktabs}     % tables
\usepackage{array}
\usepackage{xcolor}
\usepackage{tcolorbox}
\usepackage{enumitem}
\usepackage{multicol}
\usepackage{microtype}
\usepackage{parskip}
\usepackage{titlesec}
\usepackage{lmodern}

% ---------- couleurs & boîtes --------------------------------
\definecolor{myblue}{RGB}{30,100,180}
\definecolor{mygray}{RGB}{240,240,245}
\definecolor{mygreen}{RGB}{0,120,60}

\tcbuselibrary{skins,breakable}
\newtcolorbox{defbox}[1]{
  colback=mygray, colframe=myblue, fonttitle=\bfseries\small,
  title=#1, breakable, left=3pt, right=3pt, top=2pt, bottom=2pt}
\newtcolorbox{formulebox}[1]{
  colback=white, colframe=mygreen, fonttitle=\bfseries\small,
  title=#1, breakable, left=3pt, right=3pt, top=2pt, bottom=2pt}

% ---------- sections compactes -------------------------------
\titleformat{\section}{\normalfont\large\bfseries\color{myblue}}{\thesection}{0.5em}{}[\titlerule]
\titleformat{\subsection}{\normalfont\normalsize\bfseries\color{myblue}}{\thesubsection}{0.5em}{}
\titlespacing{\section}{0pt}{6pt}{3pt}
\titlespacing{\subsection}{0pt}{4pt}{2pt}

% ---------- macros -------------------------------------------
\newcommand{\R}{\mathbb{R}}
\newcommand{\E}{\mathbb{E}}
\newcommand{\Var}{\operatorname{Var}}
\newcommand{\Cov}{\operatorname{Cov}}
\newcommand{\tr}{\operatorname{tr}}
\newcommand{\diag}{\operatorname{diag}}
\newcommand{\argmax}{\operatorname*{arg\,max}}
\newcommand{\argmin}{\operatorname*{arg\,min}}
\newcommand{\T}{^{\!\top}}   % transpose

% ============================================================
\begin{document}

\twocolumn[{%
\begin{@twocolumnfalse}
  \begin{center}
    {\LARGE\bfseries\color{myblue} Synthèse MTH8304}\\[4pt]
    {\large Apprentissage non supervisé \& Séries temporelles}\\[2pt]
    {\small Polytechnique Montréal — Feuille d'examen}
  \end{center}
  \vspace{6pt}
\end{@twocolumnfalse}
}]

% ============================================================
\section{K plus proches voisins (KNN)}
% ============================================================

\begin{formulebox}{Règle de décision}
Pour classer $\mathbf{x}$, trouver les $k$ observations les plus proches selon une distance $d$ (souvent Euclidienne) et affecter la classe majoritaire :
\[
\hat{y}(\mathbf{x}) = \argmax_{c} \sum_{i \in \mathcal{N}_k(\mathbf{x})} \mathbf{1}[y_i = c]
\]
où $\mathcal{N}_k(\mathbf{x})$ est l'ensemble des $k$ voisins.
\end{formulebox}

\begin{itemize}[leftmargin=*,itemsep=1pt]
  \item Distance Euclidienne : $d(\mathbf{x},\mathbf{x'}) = \|\mathbf{x}-\mathbf{x'}\|_2$
  \item Grand $k$ $\Rightarrow$ frontière lisse (biais $\uparrow$, variance $\downarrow$)
  \item Petit $k$ $\Rightarrow$ frontière complexe (biais $\downarrow$, variance $\uparrow$)
  \item Choisir $k$ par validation croisée.
\end{itemize}

% ============================================================
\section{Régression logistique}
% ============================================================

\begin{formulebox}{Modèle}
\[
P(Y=1 \mid \mathbf{x}) = \sigma(\mathbf{x}\T\bm{\beta}) = \frac{1}{1+e^{-\mathbf{x}\T\bm{\beta}}}
\]
\textbf{Log-vraisemblance} (à maximiser) :
\[
\ell(\bm{\beta}) = \sum_{i=1}^n \bigl[y_i \log \hat{p}_i + (1-y_i)\log(1-\hat{p}_i)\bigr]
\]
\end{formulebox}

\begin{itemize}[leftmargin=*,itemsep=1pt]
  \item Pas de solution fermée $\Rightarrow$ descente de gradient / Newton-Raphson.
  \item Régularisation Ridge ($\ell_2$) ou Lasso ($\ell_1$) pour éviter le sur-ajustement.
  \item Multiclasse : one-vs-rest ou softmax.
\end{itemize}

% ============================================================
\section{Validation croisée \& Bootstrap}
% ============================================================

\subsection{Validation croisée $k$-fold}
\begin{formulebox}{Erreur CV}
\[
\widehat{\text{Err}}_{\text{CV}} = \frac{1}{n}\sum_{i=1}^n L\!\left(y_i,\, \hat{f}^{-\kappa(i)}(\mathbf{x}_i)\right)
\]
$\hat{f}^{-\kappa(i)}$ : modèle entraîné sans le fold $\kappa(i)$ contenant $i$.
\end{formulebox}

\subsection{Bootstrap}
\begin{formulebox}{Erreur .632}
Probabilité d'être absent d'un échantillon bootstrap : $(1-1/n)^n \approx e^{-1} \approx 0.368$.
\[
\widehat{\text{Err}}^{.632} = 0.368\,\overline{\text{err}} + 0.632\,\widehat{\text{Err}}^1
\]
$\overline{\text{err}}$ : erreur resubstitution ; $\widehat{\text{Err}}^1$ : erreur leave-one-out bootstrap.
\end{formulebox}

% ============================================================
\section{Analyse Discriminante (LDA / QDA)}
% ============================================================

\subsection{Hypothèses}
\[
\mathbf{X}\mid Y=k \;\sim\; \mathcal{N}(\bm{\mu}_k,\, \bm{\Sigma}_k)
\]

\begin{formulebox}{Règle de Bayes}
\[
\hat{y} = \argmax_k \;\bigl[\log\pi_k + \log p(\mathbf{x}\mid k)\bigr]
\]
\end{formulebox}

\subsection{LDA — Covariances égales $\bm{\Sigma}_k = \bm{\Sigma}$}
\begin{formulebox}{Fonction discriminante linéaire}
\[
\delta_k(\mathbf{x}) = \mathbf{x}\T\bm{\Sigma}^{-1}\bm{\mu}_k - \tfrac{1}{2}\bm{\mu}_k\T\bm{\Sigma}^{-1}\bm{\mu}_k + \log\pi_k
\]
Frontière de décision \textbf{linéaire} entre $k$ et $\ell$ :
\[
\{\mathbf{x}: \delta_k(\mathbf{x}) = \delta_\ell(\mathbf{x})\}
\]
\end{formulebox}

\subsection{QDA — Covariances différentes}
\begin{formulebox}{Fonction discriminante quadratique}
\[
\delta_k(\mathbf{x}) = -\tfrac{1}{2}\log|\bm{\Sigma}_k| -\tfrac{1}{2}(\mathbf{x}-\bm{\mu}_k)\T\bm{\Sigma}_k^{-1}(\mathbf{x}-\bm{\mu}_k) + \log\pi_k
\]
Frontière de décision \textbf{quadratique}.
\end{formulebox}

\begin{defbox}{LDA vs QDA}
\begin{tabular}{@{}lcc@{}}
\toprule
 & LDA & QDA \\
\midrule
Covariances & $\bm{\Sigma}_k = \bm{\Sigma}$ & $\bm{\Sigma}_k$ libres \\
Nb paramètres & $p(p+1)/2$ & $Kp(p+1)/2$ \\
Frontière & linéaire & quadratique \\
Biais & $\uparrow$ & $\downarrow$ \\
Variance & $\downarrow$ & $\uparrow$ \\
\bottomrule
\end{tabular}
\end{defbox}

\subsection{Estimation}
\[
\hat{\bm{\mu}}_k = \frac{1}{n_k}\sum_{i: y_i=k}\mathbf{x}_i, \quad
\hat{\bm{\Sigma}} = \frac{1}{n-K}\sum_{k}\sum_{i:y_i=k}(\mathbf{x}_i-\hat{\bm{\mu}}_k)(\mathbf{x}_i-\hat{\bm{\mu}}_k)\T
\]

% ============================================================
\section{K-moyennes (K-means)}
% ============================================================

\begin{formulebox}{Objectif}
Minimiser l'inertie intra-classe :
\[
J = \sum_{k=1}^K \sum_{i \in C_k} \|\mathbf{x}_i - \bm{\mu}_k\|^2
\]
\end{formulebox}

\begin{defbox}{Algorithme}
\begin{enumerate}[leftmargin=*,itemsep=1pt]
  \item Initialiser $K$ centres $\bm{\mu}_1,\ldots,\bm{\mu}_K$ (aléatoire ou K-means++).
  \item \textbf{Affectation} : $C_k \leftarrow \{i : k = \argmin_j\|\mathbf{x}_i - \bm{\mu}_j\|^2\}$
  \item \textbf{Mise à jour} : $\bm{\mu}_k \leftarrow \dfrac{1}{|C_k|}\sum_{i\in C_k}\mathbf{x}_i$
  \item Répéter 2–3 jusqu'à convergence.
\end{enumerate}
\end{defbox}

\begin{itemize}[leftmargin=*,itemsep=1pt]
  \item Convergence garantie mais vers un minimum local $\Rightarrow$ plusieurs initialisations.
  \item Choisir $K$ : coude de la courbe $J(K)$, silhouette, BIC.
  \item Sensible aux outliers et à l'échelle $\Rightarrow$ normaliser les données.
\end{itemize}

% ============================================================
\section{Mélange gaussien \& Algorithme EM}
% ============================================================

\begin{formulebox}{Modèle de mélange}
\[
p(\mathbf{x}) = \sum_{k=1}^K \pi_k\,\mathcal{N}(\mathbf{x};\bm{\mu}_k,\bm{\Sigma}_k)
\]
avec $\pi_k \geq 0$ et $\sum_k \pi_k = 1$.
\end{formulebox}

\begin{formulebox}{Log-vraisemblance}
\[
\ell(\bm{\theta}) = \sum_{i=1}^n \log\!\left(\sum_{k=1}^K \pi_k\,\mathcal{N}(\mathbf{x}_i;\bm{\mu}_k,\bm{\Sigma}_k)\right)
\]
\end{formulebox}

\begin{defbox}{Algorithme EM}
\textbf{E-step} — Responsabilités :
\[
r_{ik} = \frac{\pi_k\,\mathcal{N}(\mathbf{x}_i;\bm{\mu}_k,\bm{\Sigma}_k)}{\sum_j \pi_j\,\mathcal{N}(\mathbf{x}_i;\bm{\mu}_j,\bm{\Sigma}_j)}
\]
\textbf{M-step} — Mise à jour :
\[
n_k = \sum_{i=1}^n r_{ik}
\]
\[
\hat{\pi}_k = \frac{n_k}{n}, \quad
\hat{\bm{\mu}}_k = \frac{1}{n_k}\sum_i r_{ik}\,\mathbf{x}_i
\]
\[
\hat{\bm{\Sigma}}_k = \frac{1}{n_k}\sum_i r_{ik}(\mathbf{x}_i-\hat{\bm{\mu}}_k)(\mathbf{x}_i-\hat{\bm{\mu}}_k)\T
\]
\end{defbox}

\begin{defbox}{EM vs K-means}
\begin{tabular}{@{}lcc@{}}
\toprule
 & K-means & EM-GMM \\
\midrule
Affectation & dure (0/1) & douce ($r_{ik} \in [0,1]$) \\
Frontières & Voronoï (linéaire) & elliptiques \\
Paramètres & $\bm{\mu}_k$ seulement & $\pi_k,\bm{\mu}_k,\bm{\Sigma}_k$ \\
\bottomrule
\end{tabular}
\end{defbox}

% ============================================================
\section{Séries temporelles — Concepts de base}
% ============================================================

\begin{formulebox}{Stationnarité faible}
Un processus $\{X_t\}$ est stationnaire au sens large si :
\begin{align*}
\E[X_t] &= \mu \quad\text{(constante)}\\
\Cov(X_t, X_{t+h}) &= \gamma(h) \quad\text{(dépend seulement de }h\text{)}
\end{align*}
\end{formulebox}

\begin{formulebox}{Fonction d'autocovariance / autocorrélation}
\[
\gamma(h) = \Cov(X_t, X_{t+h}), \quad
\rho(h) = \frac{\gamma(h)}{\gamma(0)}
\]
\end{formulebox}

\begin{formulebox}{Autocorrélation partielle (PACF)}
$\phi_{hh}$ est le coefficient de la régression de $X_t$ sur $X_{t-1},\ldots,X_{t-h}$ après avoir contrôlé les lags intermédiaires.
\end{formulebox}

\begin{defbox}{Rendre une série stationnaire}
\begin{tabular}{@{}ll@{}}
\toprule
Problème & Transformation \\
\midrule
Tendance & Différenciation $\nabla X_t = X_t - X_{t-1}$ \\
Saisonnalité & Diff. saisonnière $\nabla_s X_t = X_t - X_{t-s}$ \\
Variance non const. & $\log X_t$ ou Box-Cox \\
\bottomrule
\end{tabular}
\end{defbox}

% ============================================================
\section{Processus ARMA}
% ============================================================

\begin{formulebox}{Modèles de base}
\textbf{AR($p$)} — Autorégressif :
\[
X_t = \phi_1 X_{t-1} + \cdots + \phi_p X_{t-p} + \varepsilon_t
\]
\textbf{MA($q$)} — Moyenne mobile :
\[
X_t = \varepsilon_t + \theta_1\varepsilon_{t-1} + \cdots + \theta_q\varepsilon_{t-q}
\]
\textbf{ARMA($p,q$)} :
\[
X_t = \sum_{i=1}^p \phi_i X_{t-i} + \varepsilon_t + \sum_{j=1}^q \theta_j \varepsilon_{t-j}
\]
où $\varepsilon_t \overset{\text{iid}}{\sim} \mathcal{N}(0,\sigma^2)$ (bruit blanc gaussien).
\end{formulebox}

\begin{formulebox}{Opérateur retard $B$}
\[
B X_t = X_{t-1}, \quad B^k X_t = X_{t-k}
\]
Polynômes caractéristiques :
\[
\phi(B) = 1 - \phi_1 B - \cdots - \phi_p B^p
\]
\[
\theta(B) = 1 + \theta_1 B + \cdots + \theta_q B^q
\]
ARMA$(p,q)$ : $\phi(B)X_t = \theta(B)\varepsilon_t$
\end{formulebox}

\begin{defbox}{Propriétés ACF/PACF}
\begin{tabular}{@{}lccc@{}}
\toprule
Modèle & ACF & PACF \\
\midrule
AR($p$) & décroît & coupure après $p$ \\
MA($q$) & coupure après $q$ & décroît \\
ARMA($p,q$) & décroît & décroît \\
\bottomrule
\end{tabular}
\end{defbox}

\subsection{Stationnarité d'un AR($p$)}
Toutes les racines du polynôme $\phi(z) = 0$ doivent être \textbf{à l'extérieur} du cercle unité ($|z| > 1$), ou bien les racines inverses à l'intérieur.

\subsection{Invertibilité d'un MA($q$)}
Toutes les racines du polynôme $\theta(z) = 0$ doivent être \textbf{à l'extérieur} du cercle unité.

% ============================================================
\section{ARIMA et SARIMA}
% ============================================================

\begin{formulebox}{ARIMA($p,d,q$)}
\[
\phi(B)(1-B)^d X_t = \theta(B)\varepsilon_t
\]
$d$ : ordre de différenciation pour obtenir la stationnarité.
\end{formulebox}

\begin{formulebox}{SARIMA($p,d,q$)$\times$($P,D,Q$)$_s$}
\[
\Phi(B^s)\,\phi(B)\,(1-B^s)^D(1-B)^d X_t = \Theta(B^s)\,\theta(B)\,\varepsilon_t
\]
\begin{tabular}{@{}ll@{}}
$\Phi(B^s) = 1 - \Phi_1 B^s - \cdots - \Phi_P B^{Ps}$ & (AR saisonnier)\\
$\Theta(B^s) = 1 + \Theta_1 B^s + \cdots + \Theta_Q B^{Qs}$ & (MA saisonnier)
\end{tabular}
\end{formulebox}

\begin{defbox}{Critères de sélection du modèle}
\[
\text{AIC} = -2\ell(\hat{\bm{\theta}}) + 2k
\]
\[
\text{AICc} = \text{AIC} + \frac{2k(k+1)}{n-k-1}
\]
\[
\text{BIC} = -2\ell(\hat{\bm{\theta}}) + k\log n
\]
$k$ = nombre de paramètres. \textbf{Minimiser} ces critères.\\
BIC pénalise plus $\Rightarrow$ modèles plus parcimonieux.
\end{defbox}

\begin{defbox}{Stratégie d'identification Box-Jenkins}
\begin{enumerate}[leftmargin=*,itemsep=1pt]
  \item Vérifier stationnarité (ADF, KPSS) ; différencier si besoin.
  \item Examiner ACF \& PACF pour suggérer $p$ et $q$.
  \item Estimer plusieurs modèles candidats.
  \item Comparer par AIC/AICc/BIC.
  \item Diagnostiquer les résidus : test de Ljung-Box, normalité, homoscédasticité.
\end{enumerate}
\end{defbox}

% ============================================================
\section{Tests de stationnarité}
% ============================================================

\begin{defbox}{Test ADF (Augmented Dickey-Fuller)}
$H_0$ : racine unitaire (non stationnaire).\\
$H_1$ : stationnaire.\\
\textbf{Rejeter $H_0$} si $p$-valeur $<$ seuil $\Rightarrow$ série stationnaire.
\end{defbox}

\begin{defbox}{Test KPSS}
$H_0$ : série stationnaire.\\
$H_1$ : racine unitaire.\\
\textbf{Ne pas rejeter $H_0$} est souhaitable.\\
\textit{Utiliser ADF et KPSS ensemble pour confirmer.}
\end{defbox}

% ============================================================
\section{Prévision avec ARIMA}
% ============================================================

\begin{formulebox}{Prévision optimale (MMSE)}
\[
\hat{X}_{n+h} = \E[X_{n+h} \mid X_n, X_{n-1}, \ldots]
\]
Propriété clé : pour un MA($q$), $\hat{X}_{n+h} = 0$ pour $h > q$.\\
Pour un AR($p$) :
\[
\hat{X}_{n+h} = \sum_{i=1}^p \hat{\phi}_i \hat{X}_{n+h-i}
\]
(en remplaçant les valeurs futures inconnues par leurs prévisions).
\end{formulebox}

\begin{formulebox}{Intervalle de prévision (95\%)}
\[
\hat{X}_{n+h} \pm 1.96\,\sigma_h
\]
où $\sigma_h^2 = \Var(X_{n+h} - \hat{X}_{n+h})$ croît avec $h$.
\end{formulebox}

% ============================================================
\section{Diagrammes ACF/PACF — Guide rapide}
% ============================================================

\begin{defbox}{Bandes de confiance à 95\%}
Sous $H_0$ : bruit blanc, les corrélations sont dans $\pm\dfrac{1.96}{\sqrt{n}}$.\\
Tout pic hors des bandes est statistiquement significatif.
\end{defbox}

\begin{defbox}{Identifier $p$ et $q$}
\begin{itemize}[leftmargin=*,itemsep=1pt]
  \item ACF coupe après lag $q$ \textbf{et} PACF décroît $\Rightarrow$ MA($q$)
  \item PACF coupe après lag $p$ \textbf{et} ACF décroît $\Rightarrow$ AR($p$)
  \item Les deux décroissent (exponentielle ou oscillation) $\Rightarrow$ ARMA
  \item Pics aux lags multiples de $s$ $\Rightarrow$ composante saisonnière d'ordre $s$
\end{itemize}
\end{defbox}

% ============================================================
\section{Résidus \& Diagnostics}
% ============================================================

\begin{formulebox}{Test de Ljung-Box}
$H_0$ : les résidus sont un bruit blanc.
\[
Q = n(n+2)\sum_{h=1}^H \frac{\hat{\rho}_\varepsilon(h)^2}{n-h} \;\overset{H_0}{\sim}\; \chi^2(H-p-q)
\]
\textbf{Rejeter $H_0$} si $Q > \chi^2_{1-\alpha}$ $\Rightarrow$ modèle inadéquat.
\end{formulebox}

Un bon modèle doit avoir :
\begin{itemize}[leftmargin=*,itemsep=1pt]
  \item Résidus $\approx$ bruit blanc (Ljung-Box non rejeté)
  \item Résidus $\approx$ normaux (QQ-plot, test Shapiro-Wilk)
  \item Variance des résidus constante (pas d'hétéroscédasticité)
\end{itemize}

% ============================================================
\section{ACP — Analyse en Composantes Principales}
% ============================================================

\begin{formulebox}{SVD / Décomposition spectrale}
\[
\bm{\Sigma} = \mathbf{V}\bm{\Lambda}\mathbf{V}\T
\]
Vecteurs propres $\mathbf{v}_j$ (composantes principales), valeurs propres $\lambda_1 \geq \cdots \geq \lambda_p \geq 0$.
\[
\text{Variance expliquée par } j = \frac{\lambda_j}{\sum_k \lambda_k}
\]
Scores : $\mathbf{z}_i = \mathbf{V}_L\T(\mathbf{x}_i - \bar{\mathbf{x}})$ où $\mathbf{V}_L$ = $L$ premières colonnes de $\mathbf{V}$.
\end{formulebox}

% ============================================================
\section{Formules utiles}
% ============================================================

\begin{formulebox}{Loi normale multivariée}
\[
\mathcal{N}(\mathbf{x};\bm{\mu},\bm{\Sigma}) = \frac{1}{(2\pi)^{p/2}|\bm{\Sigma}|^{1/2}}
\exp\!\left(-\tfrac{1}{2}(\mathbf{x}-\bm{\mu})\T\bm{\Sigma}^{-1}(\mathbf{x}-\bm{\mu})\right)
\]
Distance de Mahalanobis : $d_M^2 = (\mathbf{x}-\bm{\mu})\T\bm{\Sigma}^{-1}(\mathbf{x}-\bm{\mu})$
\end{formulebox}

\begin{formulebox}{Théorème de Bayes}
\[
P(Y=k\mid\mathbf{x}) = \frac{p(\mathbf{x}\mid Y=k)\,\pi_k}{p(\mathbf{x})} \propto p(\mathbf{x}\mid k)\,\pi_k
\]
\end{formulebox}

\begin{formulebox}{Vraisemblance gaussienne}
\[
\log p(\mathbf{x}\mid\bm{\mu},\bm{\Sigma}) = -\tfrac{p}{2}\log(2\pi) - \tfrac{1}{2}\log|\bm{\Sigma}| - \tfrac{1}{2}d_M^2
\]
\end{formulebox}

\vspace{4pt}
\begin{tcolorbox}[colback=myblue!10,colframe=myblue,fonttitle=\bfseries\small,title=Mémo rapide]
\begin{tabular}{@{}ll@{}}
\toprule
Tâche & Outil \\
\midrule
Clustering souple & EM-GMM \\
Clustering dur & K-means \\
Classification supervisée & LDA/QDA/KNN \\
Série non stationnaire & Différencier + ARIMA \\
Série saisonnière & SARIMA \\
Réduction de dimension & ACP \\
\bottomrule
\end{tabular}
\end{tcolorbox}

\end{document}
